\section{Discussion}
\label{sec:discussion}
The finite-sample analysis connects same-sample factor learning to estimation of the original sparse AFT slope. The reference representation preserves the ordinary-model predictor exactly, placing the effect of factor estimation in the design geometry. Conditional score centering, column-scale empirical-process control, and local event information then yield coefficient bounds and finite attainment under explicit assumptions. Additional factor effects enter through a separate approximation term relative to the specified latent representation.

This separation is also useful computationally. In the fixed strong-dependence illustrations, FA has lower maximum active curvature and smaller active spectral ratios, together with larger scalar steps and less work to common accuracy. Late UA residuals concentrate in weak-curvature directions on the selected stable supports. A single active-face Newton correction supplies the direction-specific scaling needed to complete each of those fits. The numerical policy combines proximal search with qualified local refinement, using the same full-coordinate accuracy criterion for both estimators.

Training selection gives complementary information about the representation. Some FA fits retain a dense factor-mediated molecular score with a zero residual coefficient block. Thus $\widehat{\bm{\beta}}$ and $\widehat{\bm{b}}_{\rm pred}$ answer different questions: the former describes residual associations, while the latter records the molecular slope deployed for prediction. The six training partitions exhibit different selected-score orderings across methods, making both coefficient recovery and independent prediction relevant to evaluation of the complete selection procedure.

The theoretical bounds concern fixed or covariate-measurable tuning, and the numerical experiments describe fixed-sample geometry and training selection. Independent repeated samples can assess the risk of the outcome-selected procedure under prespecified grids and numerical rules. A biomedical application would additionally require a common endpoint, an explicit observation scheme, and evaluation of the frozen prediction map on an external cohort. These evaluations extend the present theory and computational illustrations to statistical performance and applied use.

The resulting framework links three objects: the original sparse AFT target, covariate geometry learned from the sample, and the numerical accuracy of the fitted estimator. Exact reference coordinates preserve the target, finite-sample analysis accounts for censoring and factor estimation, and local curvature completes the computational connection. Together, they provide a coherent basis for sparse censored regression with estimated factor structure.
